\documentclass[12pt]{article}
\usepackage[a4paper,margin=1in]{geometry}
\usepackage{setspace}
\usepackage{titlesec}
\usepackage{xcolor}
\usepackage{changepage} % for adjustwidth
\usepackage{xparse}

% --- Bold Colored Heading ---
\NewDocumentCommand{\BC}{O{black} m}{%
  \vspace{\baselineskip}%
  \noindent\textbf{\color{#1}#2}\par
}

% --- Response Block ---
\NewDocumentCommand{\RE}{O{black}}{%
  \vspace{\baselineskip}%
  \noindent\textbf{\color{#1}{Response}:}\par
}

% --- Quotation Block ---
\NewDocumentCommand{\QT}{O{blue} +m}{%
  \begin{adjustwidth}{3em}{0pt}%
  \color{#1}\itshape
  ``#2''%
  \end{adjustwidth}%
  \par\vspace{\baselineskip}%
}

\titleformat{\section}{\large\bfseries}{\thesection}{1em}{}
\doublespacing

\begin{document}

\begin{center}
    {\Large \textbf{Response to Reviewers}}\\[1em]
    Manuscript ID: 25-06155\\
    Title: Climate-driven scour impacts on bridge capacities: a distribution-agnostic framework
\end{center}


\doublespacing

\section*{General Note}
\RE We thank the editors and reviewers for their careful reading of the manuscript and their constructive comments and suggestions. We have revised the paper accordingly, incorporating significant revisions (including newly added statements, paragraphs, and citations, as well as shortened presentations, as requested by the reviewers). In addition, we have thoroughly proofread the manuscript and corrected all grammatical errors.

In the following, the original comments and suggestions from the Editor and the reviewers are shown in {\color{black}{\textbf{bold fonts}}} in chronological order, followed by our point-by-point responses; if revision is significant, we will quote the revised text in {\color{blue}{\textit{blue italic fonts}}}.

\newpage

\section*{Reviewer \#1}

\BC{The surrogate model is trained using only a single input parameter (scour depth) with a tuple of output parameters. While feature engineering was attempted (e.g., $log(Sz)$, $1/Sz$, and $\sqrt{Sz}$), this is not sufficiently robust or generalisable. It is recommended to include additional input variables, such as material properties, span length, pier height, and other structural or geotechnical parameters, treated as random variables from appropriate distributions. This would improve model generality and applicability.}


\RE{That's a fair point. Thanks for pointing it out and we agree with your general remark. However, please allow us to clarify. 

\RE{We appreciate this constructive suggestion and agree with the reviewer’s general remark. 

\begin{itemize}
    \item Our finite element (FE) simulations did incorporate material uncertainties—specifically, variations in concrete and steel strengths—as random variables. However, for the surrogate modeling stage, we intentionally used only the scour depth as the explicit input variable. This design choice was made to isolate and systematically analyze the causal influence of scour, which is the central focus of this study. Regarding other variables, such as geometric dimensions, their variations are more insignificant. Please see the quoted text below. 

    \item The fundamental reason for excluding other inputs in the surrogate model is that the fluctuations in material properties (e.g., concrete and steel strength) were non-causal with respect to scour-induced capacity degradation. These uncertainties were therefore treated as background noise within the simulation dataset, allowing the surrogate to learn the direct, causal mapping between scour depth and bridge capacity response. 
    \item We acknowledge, however, that if a broader climate-change context is considered—where increasing temperature, humidity, or corrosion may causally affect material degradation—then extending the surrogate framework to a multi-input formulation becomes essential. This has been highlighted as a key direction for future work in the revised Discussion section. Please see the revised Dicussion quoted below. 
\end{itemize}
}

Regarding uncertainties in other material variables: 
\QT{`Specifically, concrete compressive strength $f'_c$ is modeled as a normal distribution with a mean of 27.0 MPa and a standard deviation of 1.0 MPa. Steel yield strength $f_y$ is modeled using a lognormal distribution with a mean of 420.0 MPa and a standard deviation of 4.2 MPa. Latin Hypercube Sampling (LHS) is employed with 1,000 samples per scenario (three scenarios in total), resulting in 3,000 simulations. This sampling approach robustly approximates the probabilistic scour hazard while keeping material variability constrained ... Figure 5a and 5b depict the resulting histograms based on the sampled values for concrete and reinforcing steel strength, respectively.'}

In addition, we have improved the first disucssion item:
\QT{`In this work, scour is treated as the dominant causal source of uncertainty when assessing bridge capacity degradation under climate scenarios. Material uncertainties (i.e., concrete compressive strength and steel yield strength) were incorporated in the FE simulations as random variables but were treated as background fluctuations rather than causal drivers of scour effects. 
We recognize that, in a more comprehensive climate-change context, variables such as ambient temperature, moisture, and corrosion may lead to progressive material degradation and thus partially and causally affect bridge capacity. In such cases, a multi-input surrogate modeling framework that integrates both environmental and material factors would become essential for accurate and generalizable prediction.'}



\BC{The hyperparameter tuning process for SVR and GBR is mentioned but not detailed (e.g., parameter search ranges, cross-validation folds, optimisation criteria). This information is essential for reproducibility.}

\RE{We appreciate the reviewer’s emphasis on reproducibility. The key hyperparameters that were actively tuned—such as the penalty parameter $C$ for SVR, and the learning rate and maximum tree depth for GBR—are listed in Table 2. For parameters not shown in that table, the standard default values from the \texttt{scikit-learn} library were adopted. 

In the revised manuscript, we have clarified that hyperparameter selection was performed using a grid search combined with 5-fold cross-validation, optimising the mean squared error (MSE) criterion on the validation folds. This ensures that model tuning was both systematic and transparent. These additions now make the process fully reproducible. \textit{We have added the following new text to justify our treatment:
}}


\QT{`Hyperparameter selection for both SVR and GBR was conducted via grid search with 5-fold cross-validation using bootstrapped datasets. Specifically, the optimisation objective was to minimise the mean squared error (MSE) across validation folds. For SVR, the penalty parameter $C$ and kernel parameter $\gamma$ were varied within logarithmic ranges $[10^{-2}, 10^{3}]$, while for GBR, the learning rate and maximum tree depth were tuned within $[0.01, 0.2]$ and $[2, 6]$, respectively. Other parameters followed the default settings of the \texttt{scikit-learn} implementation. The final hyperparameter settings are provided in Table 2' 
}

\BC{The FE model is based on a single bridge configuration. A discussion is needed on how the surrogate model's predictive capability extends to other bridge typologies or geometric configurations.}
\RE{As it stands, the trained surrogate model is specific to the single bridge we analyzed—a typical Midwestern U.S. bridge with four spans. The model in this paper is best seen as a detailed case study demonstrating the framework: linking nonlinear FE simulations to a fast surrogate model and then applying our credal-set uncertainty analysis. To make this tool predictive for other bridge types, we would need to significantly expand our initial FE simulation dataset. Instead of just sampling scour depth and material properties, we would also have to vary key geometric and structural parameters—like pier height, span length, and foundation type—as inputs. So, while this specific model isn't transferable, the methodology is. It's adaptable to any bridge type, but it would require generating a new, broader set of training data specific to those designs.}

\BC{The study only evaluates transverse (lateral) capacity under quasi-static pushover analysis. Real flood-induced loadings can be dynamic and multi-directional. A discussion on extending the framework to dynamic loading cases and/or multi-hazard coupling (e.g., scour + seismic) is warranted.}
\RE{We focused on static pushover analysis as it's a standard method to define a bridge's fundamental capacity, allowing us to first model how scour degrades that capacity. However, this framework is flexible. The pushover simulation could be replaced with dynamic analyses to model time-varying loads. Furthermore, as noted in our literature review and discussion, the surrogate model could be extended to handle multi-hazard scenarios by simply adding more inputs, such as seismic intensity, alongside scour depth.}

\BC{It is unclear why only SVR and GBR were chosen. Ideally, these should be selected from a broader set of potential methods, with a comparative performance analysis demonstrating their superiority. Any limitations of the chosen methods should also be explicitly stated.}
\RE{We selected SVR and GBR because they are robust, widely-used in structural engineering, and have complementary strengths. SVR is excellent for smooth, nonlinear data, while GBR is very effective at capturing the sharp, non-smooth transitions we saw in our results. Benchmarking studies show they perform well without the complexity of deep learning. We also discussed their limitations, such as SVR's tendency to smooth over the abrupt “jumps” that GBR captured.}

\section*{Reviewer \#2}

\BC{In the Introduction section, could the author give the definition of “scour-critical bridge”?}
\RE{A scour-critical bridge refers to a bridge whose foundation is unstable due to observed scour or is predicted to become unstable during design flood events, as defined by the U.S. Federal Highway Administration (FHWA, 2012).}

\BC{Three scenarios (Moderate, High, Extreme) of flood velocities are considered in this manuscript. However, Scenario 1 and 2 are from Missouri River and Colorado River respectively. Why did the authors select data from two different rivers?}
\RE{The Missouri River and Colorado River were selected as spatial analogues because their flood regimes capture a representative spectrum of climate-driven extremes relevant to Midwestern bridges. The Missouri River reflects a moderate-to-high flow regime typical of large Midwestern basins, whereas the Colorado River provides an analogue for extreme, high-velocity flood scenarios. This analogue-based approach was adopted to bracket the uncertainty space, rather than relying on an “average river,” which may underestimate the potential impact of rare but critical extremes.}

\BC{How do the authors consider traffic load and crowd load in FE modeling?}
\RE{For simplicity, live loads such as vehicular traffic and crowd loading were not explicitly included in the finite element analyses. This choice was made to isolate the effect of scour depth on structural capacity and to highlight the surrogate modeling framework. In practice, live loads can be incorporated as additional stochastic variables, which will be considered in future extensions of this framework.}

\BC{In section 3.5, why do the authors select these features? Other features seem to have no physical significance.}
\RE{These transformations were introduced to enrich the input space and capture nonlinear patterns observed in the finite element outputs. While not all transformations have direct physical interpretation, they improve the surrogate’s ability to approximate the nonlinear relationship between scour depth and structural capacity.}

\BC{Please discuss the selection of hyperparameters of surrogate model.}
\RE{The hyperparameters of SVR (kernel type, C, $\epsilon$, $\gamma$) and GBR (number of estimators, learning rate, maximum depth) were optimized through a grid search with 5-fold cross-validation. The search ranges were selected based on prior studies and scikit-learn guidelines. The final values reported in Table 2 correspond to the combination that minimized the cross-validated RMSE while avoiding overfitting. We have clarified this procedure in Section 3.5 and added a supplementary figure illustrating the sensitivity of model accuracy to key hyperparameters.}

\BC{In Figure 8 and Figure 9, why is there a sudden rise in the raw simulation data, and why can’t the surrogate model reflect it?}
\RE{The localized “rise” in capacity observed around \(S_z \approx  8 - 10\ \) m in the raw FE simulations is attributed to nonlinear soil–structure interaction effects and geometric/material nonlinearities that temporarily redistribute demand. This branching behavior is not sustained and is followed by rapid degradation at larger scour depths. Surrogate models (SVR, GBR) approximate the dominant monotonic trend but cannot reproduce these abrupt irregularities, which represent secondary nonlinear effects rather than the primary capacity–scour relationship.}

\section*{Reviewer \#3}

\BC{While the paper claims to be “distribution-agnostic” and to introduce a “credal-set” approach, the novelty over existing uncertainty quantification methods needs to be emphasized more clearly.}
\RE{We appreciate this point. Our credal-set approach differs fundamentally from Bayesian or conformal methods in that it does not assume any underlying probability distribution for input uncertainties. Instead, it bounds the uncertainty space through interval-valued probability sets, providing robustness to deep uncertainty. We will clarify these distinctions and add a conceptual diagram contrasting our approach with probabilistic UQ in the revised manuscript.}

\BC{The surrogate models are trained entirely on simulated data from OpenSeesPy models. External validation would enhance credibility.}
\RE{The surrogate model is only as good as the simulation data it was trained on, and it currently lacks external validation against field or experimental data. Our paper focuses on presenting the computational framework, and we agree that this real-world validation is the necessary next step. Because we used only one bridge configuration and one set of soil-structure assumptions, our specific surrogate model is not generalizable. Any biases or limitations in our original OpenSeesPy model are inherently learned and propagated by the surrogate.}

\BC{The “spatial analogue” method for scenario selection may introduce its own uncertainties. Discuss the representativeness and calibration of the chosen rivers.}
\RE{We chose the river velocities not to imply a perfect transfer, but to bracket the range of plausible uncertainties. The Missouri River represents a typical average flood, while the Colorado represents a peak flood. We also included a “worst-case” 500-year flood to ensure full coverage of conditions. The scour profiles were not calibrated to a specific historical event; instead, flow velocities were used as inputs to the standard HEC-18 erosion framework, which computed scour depths using established equations.}

\BC{Table 2 lists hyperparameters, but it is unclear how sensitive the results are to these parameter choices.}
\RE{The hyperparameters in Table 4 (SVR) and Table 5 (GBR) were selected to balance model accuracy with generalization and avoid overfitting. The GBR model used many shallow trees (max\_depth=3) with a small learning rate (0.015) for stability, and the SVR used a high penalty (C=100) and tight margin (\(\xi=0.01\)) to closely fit data. Engineered features were chosen manually to reflect expected scour-response physics, such as exponential or asymptotic decay. No formal multicollinearity test was done; instead, feature dominance was assessed post-training (Figures 14–15).}

\BC{Figures 8 and 9 show prediction intervals, but their interpretation could be expanded.}
\RE{The intervals widen significantly in regions of sparse training data, such as at extreme scour depths (> 12.5 m), which intentionally indicates reduced model confidence. They also widen in dense regions where raw data is highly variable, e.g., around \(S_z \approx 7.5 \) m, due to modeling artifacts. Table 7 lists bounds at key depths, but we agree additional quantitative measures such as average interval width or coverage probability would enhance the analysis.}

\BC{The conclusion states that the framework supports resilient decision-making but lacks practical guidance.}
\RE{Imagine a transportation agency with a safety threshold of 8,000 kN for base shear. For a 15 m scour scenario, a single simulation might predict 8,500 kN, appearing “safe.” Our framework, however, provides a credal-set interval: median 8,489 kN and lower bound 7,346 kN. Using the lower bound, the agency correctly flags the bridge as high-priority for retrofit—illustrating how credal-set analysis supports conservative, resilience-oriented decision-making.}

\section*{Reviewer \#4}

\BC{The paper suffers from organizational issues in its content structure. It devotes excessive space to background while lacking methodological detail.}
\RE{We appreciate the reviewer’s observation. In the revised version, we have restructured Section 3 to emphasize methodology and validation, reducing redundancy in background material and improving logical flow between model development, surrogate training, and uncertainty quantification.}

\BC{Figure 2: The flowchart is not expressed clearly enough and fails to present the complete logical chain for conducting the research.}
\RE{We apologize for the lack of clarity. The logical chain is: (1) generate 3,000 probabilistic inputs for scour and material properties; (2) build 3,000 FE bridge models in OpenSeesPy; (3) conduct nonlinear pushover analyses; (4) extract the four-part “capacity tuple”; (5) train SVR and GBR surrogate models; (6) evaluate accuracy and apply bootstrap ensembles for credal-set bounds. This will be clarified and visually improved in the revised Figure 2.}

\BC{The applicability of Formula (2) and Formula (3) needs further verification regarding Reynolds number assumptions.}
\RE{We acknowledge this limitation. In the revision, we will clarify that these empirical relations are intended as dimensionally consistent parameterizations within the valid Reynolds regime for bridge-scale flow, and not as universal predictive formulas. The issue of negative scour depths at extreme Reynolds numbers will be explicitly addressed.}

\BC{Hydrodynamic loads at high flow velocities were not mentioned.}
\RE{Our goal was to isolate structural capacity degradation due to scour. The hydrodynamic thrust during flood events indeed imposes additional load demands that should be considered in a comprehensive vulnerability assessment. We acknowledge this omission and will discuss it as a limitation, noting that multi-hazard coupling (scour + flow + seismic) is a logical future extension.}

\BC{Material variability was minimized to highlight scour impact, but in reality, failures involve multiple coupled factors.}
\RE{That’s an excellent point. By minimizing material variability (tight standard deviations for \(f'_c\) and \(f_y\)), the surrogate model learned the scour–capacity link but not coupled effects. In practice, weaker materials may exacerbate scour-induced degradation. We will expand our discussion to acknowledge that the framework currently underestimates total uncertainty and that future models should include material properties as explicit stochastic inputs.}

\BC{No evidence for model validation is provided.}
\RE{We agree. Our FE model was based on Atak et al. (2014), but we did not include a direct validation against that study. We recognize this as a gap and will include comparison results or reference validation steps in the revision, ensuring consistency with prior literature.}

\BC{Polynomial and logarithmic transformations lack clear physical support.}
\RE{These transformations were chosen to reflect expected nonlinear scour–capacity behavior (e.g., power-law or exponential decay). Polynomial (\(S_z^2, S_z^3\)) and logarithmic (\(\log S_z\)) terms capture curvature, while inverse and square-root terms model asymptotic trends such as diminishing returns at deep scour depths. We will make this physical rationale more explicit.}

\BC{GBR parameters are fixed; provide tuning justification.}
\RE{We used shallow trees (max\_depth = 3) and a small learning rate (0.015) to ensure stable convergence and avoid overfitting. While parameter tuning was performed heuristically, adding a learning curve or more formal cross-validation results would improve transparency, which we plan to include.}

\BC{Clarify training/test split and representation of deep scour cases.}
\RE{The held-out test set was generated via random 80/20 partitioning. We agree that this may underrepresent extreme scour depths; hence, our credal-set intervals explicitly widen in those sparse regions to reflect model uncertainty. We will clarify this data-splitting strategy in Section 4.1.}

\BC{Global error metrics may mask local performance differences.}
\RE{Excellent observation. While global \(R^2\) and RMSE summarize overall fit, local errors in deep scour zones are better captured by our credal-set bounds. Table 7 lists bounds at specific depths (10 m, 15 m), directly revealing prediction spread. This localized uncertainty quantification is central to our approach and mitigates the masking effect of global metrics.}

\end{document}
